\chapter{Historical Background, Key Challenges and Health Outcomes}

Healthcare System in Uzbekistan

\section{Historical Background and Current Organization}\label{historical-background-and-current-organization}

Uzbekistan's health system is a product of its Soviet legacy, shaped by
the Semashko model---a highly centralized, state-funded, and
state-delivered system emphasizing hospital-based care, vertical disease
programs, and quantitative targets for inputs such as beds and staff
(Bernd Rechel M. A., 2014). During the Soviet era, health care was
nominally free at the point of use, with a focus on communicable disease
control, maternal and child health, and the expansion of physical
infrastructure. However, the system was characterized by inefficiencies,
underfunding, and a neglect of primary care and quality assurance (Bernd
Rechel E. R., 2014).

Following independence in 1991, Uzbekistan faced the dual challenge of
economic transition and health system reform. The early 1990s saw a
sharp contraction in public spending, a decline in health indicators,
and the emergence of informal payments and private provision (Bernd
Rechel M. A., 2014). The government responded with a series of reforms
aimed at rationalizing hospital capacity, strengthening PHC, and
introducing new financing mechanisms.

\textbf{Key Reform Milestones:}

\begin{itemize}
\item
  \textbf{1996 Law on Health Protection:} Established the legal basis
  for a state-guaranteed benefits package, including primary care,
  emergency care, and care for ``socially significant and hazardous''
  conditions (Bernd Rechel M. A., 2014).
\item
  \textbf{1998 Presidential Decree:} Launched a state program for health
  system reform, prioritizing PHC, emergency care, and the development
  of the private sector (Bernd Rechel M. A., 2014).
\item
  \textbf{2000s--2010s:} Implementation of World Bank and ADB-supported
  projects (Health I, II, III; Woman and Child Health Development)
  focused on PHC infrastructure, maternal and child health, and
  financing reforms (Bernd Rechel M. A., 2014). (The World Bank, 2025)
\item
  \textbf{2018--2021:} Adoption of the ``Concept on Health Development
  2019--2025'' (Presidential Decree 5590), establishment of the SHIF,
  and piloting of strategic purchasing and digital health initiatives in
  Syrdarya (The World Bank, 2025).
\end{itemize}

\textbf{Current Organization:}

The health system is organized into three administrative tiers:

\begin{itemize}
\item
  \textbf{National (Republican) Level:} Ministry of Health, national
  research institutes, tertiary hospitals, and regulatory agencies.
\item
  \textbf{Regional (Viloyat) Level:} Regional health authorities,
  multi-specialty hospitals, and specialized centers.
\item
  \textbf{District/City Level:} District medical unions, central
  hospitals, family polyclinics, and rural physician points (Bernd
  Rechel M. A., 2014).
\end{itemize}

The MoH remains the principal steward, with direct management of
national-level institutions and regulatory oversight of regional and
district facilities. The private sector, while growing, is still limited
in scope and subject to strict licensing and regulatory controls (The
World Bank, 2025).

\textbf{Service Delivery Model:}

\begin{itemize}
\item
  \textbf{Primary Care:} Delivered through family doctor points,
  polyclinics, and rural physician posts. PHC reforms have aimed to
  shift from specialist-dominated polyclinics to generalist-led family
  medicine, but progress has been uneven (Bernd Rechel A. S., 2023).
\item
  \textbf{Secondary and Tertiary Care:} Provided by district, regional,
  and national hospitals, with a legacy of overcapacity and
  fragmentation. Recent reforms aim to consolidate specialized centers
  into multi-specialty general hospitals (The World Bank, 2025)
\item
  \textbf{Emergency Care:} A well-developed network of emergency
  departments often better equipped than other public facilities, but
  subject to overload due to gaps in PHC and benefits coverage (World
  Health Organization, 2022).
\end{itemize}

\textbf{Governance and Regulation:}

Regulation is highly centralized, with the MoH and Cabinet of Ministers
responsible for policy, financing, and oversight. Local authorities have
limited autonomy, and dual accountability structures persist (Bernd
Rechel M. A., 2014). The regulatory framework for private providers and
pharmaceuticals is evolving, with recent efforts to strengthen
licensing, quality assurance, and market surveillance (The World Bank,
2025).

\section{Key Challenges in Access, Quality, and Financing}\label{key-challenges-in-access-quality-and-financing}

\textbf{Access:}

Despite improvements in health indicators, access to quality care
remains uneven. Rural-urban disparities are pronounced, with rural areas
facing shortages of qualified staff, inadequate infrastructure, and
limited access to diagnostics and medicines (Bernd Rechel A. S., 2023).
For example, only 57\% of PHC facilities reported basic water services
and 26\% basic sanitation in 2020, reflecting broader readiness
constraints (World Health Organization, 2022).

Geographical access is further constrained by the distribution of
facilities and workforce. While the number of physicians per 100,000
population has increased to 276 by 2019, and nurse density to
\textasciitilde1,100 per 100,000, these resources are concentrated in
urban centers. Outward migration of health professionals, especially to
Russia and Kazakhstan, exacerbates shortages in underserved areas.
(Bernd Rechel E. R., 2014)

\textbf{Quality:}

Quality of care is undermined by several factors:

\begin{itemize}
\item
  \textbf{Legacy of Input-Based Planning:} The Soviet system prioritized
  quantitative targets (beds, staff) over outcomes, with limited
  attention to quality assurance, patient safety, or evidence-based
  practice (Ketevan Glonti, 2013).
\item
  \textbf{Outdated Clinical Protocols:} Many facilities rely on outdated
  or poorly adapted clinical guidelines, with limited adoption of
  international best practices (Bernd Rechel M. A., 2014).
\item
  \textbf{Weak Monitoring and Evaluation:} Data systems are fragmented,
  paper-based, and focused on structural indicators rather than process
  or outcome measures. Quality improvement initiatives are limited in
  scope and sustainability (The World Bank, 2025)
\item
  \textbf{Patient Experience:} User experience is often neglected, with
  long waiting times, limited appointment systems, and a lack of
  patient-centered care models (Bernd Rechel M. A., 2014).
\end{itemize}

\textbf{Financing:}

The health financing system is characterized by:

\begin{itemize}
\item
  \textbf{Low Public Spending:} Public spending on health was
  \textasciitilde2.3\% of GDP in 2019, above the Central Asia average
  but well below the WHO European Region average (\textasciitilde5\%)
  (World Health Organization, 2022).
\item
  \textbf{High Out-of-Pocket Payments:} OOP payments accounted for
  57.7\% of current health expenditure in 2019, exposing households to
  financial hardship and catastrophic health spending. In 2020, 18\% of
  households reported cost-related nonuse of care (World Health
  Organization, 2022)
\item
  \textbf{Fragmented Pooling and Purchasing:} Budgeting is historically
  incremental, with limited strategic purchasing or risk pooling. The
  SHIF pilot in Syrdarya aims to introduce capitation and DRG-based
  payments, but implementation is nascent and faces capacity constraints
  (The World Bank, 2025).
\item
  \textbf{Limited Benefits Coverage:} The basic benefits package covers
  primary and emergency care, but excludes much of secondary/tertiary
  care and outpatient medicines, amplifying OOP and incentivizing
  emergency use (The World Bank, 2025).
\end{itemize}

\textbf{Equity and Financial Protection} (Bernd Rechel M. A.,
2014)\textbf{:}

\begin{itemize}
\item
  \textbf{Vulnerable Groups:} While the benefits package includes
  provisions for vulnerable groups (e.g., children, pregnant women,
  people with disabilities), eligibility criteria are not directly
  linked to income, and coverage gaps persist.
\item
  \textbf{Informal Payments:} Informal payments remain widespread,
  particularly in secondary and tertiary care, undermining equity and
  trust in the system.
\end{itemize}

\section{Health Outcomes and Risk Profile}\label{health-outcomes-and-risk-profile}

\textbf{Demographic and Epidemiological Trends} (WHO Regional Office for
Europe, 2023)\textbf{:}

\begin{itemize}
\item
  \textbf{Population:} 33.3 million in 2020, with a young age structure
  but a growing burden of NCDs.
\item
  \textbf{Life Expectancy:} Increased from 67.2 years in 2000 to 73.9
  years in 2016, with a gender gap of 4.7 years.
\item
  \textbf{Maternal and Infant Mortality:} Maternal mortality declined
  from 41 per 100,000 live births in 2000 to 29 in 2017; infant
  mortality from 51.8 per 1,000 in 2000 to 15.6 in 2019.
\item
  \textbf{NCDs:} Account for 84\% of all deaths in 2018, with ischemic
  heart disease and stroke as leading causes.
\item
  \textbf{Risk Factors:} High systolic blood pressure (30.7\%), dietary
  risks (28.2\%), and high fasting plasma glucose (20.2\%) are dominant
  contributors to mortality.
\end{itemize}

\textbf{Service Coverage and UHC Index} (World Health Organization,
2022)\textbf{:}

\begin{itemize}
\item
  \textbf{UHC Service Coverage Index:} Improved from 56 (2000) to 71
  (2019), below OECD/EU comparators but signaling progress.
\item
  \textbf{Immunization:} High coverage rates for childhood vaccines, but
  challenges remain in TB, HIV/AIDS, and MDR-TB.
\end{itemize}

\textbf{Social Determinants} (The World Bank Group, 2021)\textbf{:}

\begin{itemize}
\item
  \textbf{Poverty:} Poverty headcount ratio declined from 14.1\% in 2013
  to 11.4\% in 2018, but disparities persist.
\item
  \textbf{Water and Sanitation:} Infrastructure needs extensive
  rehabilitation, with poor access in many areas.
\end{itemize}

